\begin{center}
\section*{R}
\end{center}
\addcontentsline{toc}{section}{R}
\term{Repository}
Uno spazio condiviso tra i membri del gruppo, in cui viene memorizzato e versionato il codice sorgente prodotto.
\term{Requisiti di vincolo}
Requisito che descrive le restrizioni imposte dall'azienda proponente che il progetto deve rispettare.
\term{Requisito funzionale}
Definiscono cosa il sistema deve fare, descrivendo le azioni e funzioni specifiche che deve offrire.
\term{Requisito non funzionale}
Specifica come il sistema deve comportarsi mentre svolge quelle funzioni, si concentrano su qualità come le prestazioni, la sicurezza o l'affidabilità.
\term{Responsabile}
Figura che coordina le attività del gruppo e garantisce una pianificazione efficace.
\term{REST}
Representational State Transfer, è uno stile architetturale per lo sviluppo di servizi web che utilizza protocolli standard (solitamente HTTP) per operazioni su risorse identificate tramite URL, basandosi su metodi come GET, POST, PUT, DELETE.
\term{RTB}
Baseline iniziale e vincolante di requisiti e tecnologie approvate per il progetto. Comprende requisiti funzionali e non funzionali, insieme a strumenti, linguaggi e infrastrutture scelte. Costituisce il riferimento stabile per verificare coerenza e avanzamento.
